\begin{frame}{Motivation}
  \begin{itemize}
  \item Classification involving \emph{high-dimensional} and \emph{imbalanced} datasets is challenging
  \item Often encountered in domains such as:
    \begin{itemize}
    \item Medicine \& Healthcare
    \item Genomic Research % Genomforschung
    \item Drug Discovery   % Arzneimittelforschung
    \end{itemize}
  \item Especially in applications involving high-risk decision-making, AI has still been short on delivery \cite{berisha2021}
  % Insbesondere bei Anwendungen, bei denen risikoreiche Entscheidungen getroffen werden müssen, hat die KI bislang noch nicht die Erwartungen erfüllt
  % Für eine ausführliche Behandlung des Themas kann ich das Paper hier empfehlen
    \begin{itemize}
    \item Curse of dimensionality: Few samples but high number of features \cite{berisha2021}
    \item High degree of class imbalances (?)
    \end{itemize}
  \item \textbf{Feature selection:}\\Reduce dimensionality by removing irrelevant and/or redundant features
  \end{itemize}

  \begin{textblock*}{2.2cm}(13.2cm, 5.75cm)
    \centering
    \begin{tabular}{r}
      {\tiny\color{udeBlue}\textbf{On challenges in}} \\ [-0.25cm]
      {\tiny\color{udeBlue}\textbf{high-risk applications:}} \\
      \qrcode[height=1.5cm]{https://doi.org/10.1038/s41746-021-00521-5}
    \end{tabular}
  \end{textblock*}
\end{frame}

\begin{frame}{Motivation -- Regarding Class Imbalance}
  \citet{prati2004} showed that very often class imbalance has been unfairly blamed for bad performance, when (quite) often \emph{class overlapping} is a larger culprit
  \begin{columns}[T]
    \begin{column}{0.52\textwidth}
      \hspace{0.05em} Take the original dummy example used in \cite{prati2004}:
      \begin{itemize}
      \item Positive priors vary (i.e. ``height'' of the distribution curve)
      \item Solid vertical line denotes optimal Bayes split
      \item Distributions are Gaussian and both have the same standard deviation (SD) but:
        \begin{itemize}
        \item (a) Means are 1 SD apart
        \item (b) Means are 4 SD apart
        \end{itemize}
      \end{itemize}
    \end{column}
    \begin{column}{0.48\textwidth}
      $\quad$ \\ [0.5cm]
      \footnotesize Solid -- Negative Class | Dotted -- Positive Class
      \\ [0.5cm]
      \resizebox{\linewidth}{!}{\input{PratiPicture.pdf_tex}}
    \end{column}
    \end{columns}
\end{frame}

\begin{frame}{Motivation -- Human-in-the-Loop Machine Learning}
  \begin{columns}[T]
    \begin{column}{0.52\textwidth}
      \begin{itemize}
      \item Pure data-driven techniques are inadequate (especially for high-risk decision-making) \cite{maadi2021}
        \begin{itemize}
          \item Quality highly depends on input data!
        \end{itemize}
      \item \textbf{Human-in-the-Loop Machine Learning:}\\Human experts intervene when models fail to generalize
        \begin{itemize}
        \item One key factor for this approach is \emph{transparency}
        \end{itemize}
      \end{itemize}
    \end{column}
    \begin{column}{0.48\textwidth}
    \end{column}
  \end{columns}
\end{frame}

\begin{frame}{Why Another Filter Feature-Selection Method?}
  \begin{columns}[T]
    \begin{column}{0.52\textwidth}
      \textbf{Limitations of Existing Approaches:}
      \begin{itemize}
        \item \textbf{Opaque Metrics:} Mutual Information (bits/nats) or ReliefF weights offer little geometric intuition.
        \item \textbf{Inductive Bias:} Hybrid methods (e.g., FeSNOC) couple feature ranking to downstream classifiers ($k$NN).
        \item \textbf{Memory Bottlenecks:} Instance-based filters scale at $\mathcal{O}(N^2)$ memory for distance matrices.
      \end{itemize}
    \end{column}
    \begin{column}{0.45\textwidth}
      \begin{block}{Core Idea of ProD}
        Evaluate each feature independently as a \textbf{pure filter} by directly measuring the geometric \textbf{Niche Overlapping Coefficient (NOC)} of its class-wise Kernel Density Estimates (KDE).
      \end{block}
      \vspace{0.3em}
      \begin{itemize}
        \item Bounded, intuitive score in $[0, 1]$
        \item Directly visualizable for human-in-the-loop diagnostics
      \end{itemize}
    \end{column}
  \end{columns}
\end{frame}
